\section{Example Outputs}
\label{app:examples}

We present representative article excerpts from each condition for the 2024 U.S.\ Presidential Election event (\$3.69B trading volume, actual outcome: Donald Trump).

\para{\nomarket.}
\begin{quote}
\small
\emph{``Joe Biden Wins Re-Election in Historic 2024 Presidential Race. In a closely contested election that captivated the nation, President Joe Biden has secured a second term in the White House, defeating Republican challenger Donald Trump in what political analysts are calling a referendum on American democracy\ldots''}
\end{quote}

\para{\directmarket.}
\begin{quote}
\small
\emph{``Donald Trump Wins 2024 Presidential Election. Republican candidate Donald Trump has won the 2024 United States Presidential Election, according to results confirmed across key battleground states. The \$3.69 billion in trading volume on major prediction markets had signaled Trump's victory in the final weeks\ldots''}
\end{quote}

\para{\mcp.}
\begin{quote}
\small
\emph{``Donald Trump Clinches 2024 Presidential Election in Stunning Comeback. Donald Trump has secured the presidency in a result that aligns with robust prediction market signals observed throughout the campaign's final months. Trump's win, confirmed by official tallies from key battleground states, reflects shifting voter priorities around economic concerns and immigration policy\ldots''}
\end{quote}

\para{\superforecaster.}
\begin{quote}
\small
\emph{``Analysis: How Trump Defied the Polls to Win 2024. The 2024 presidential election underscores a pattern that superforecasters have long recognized: base rates and structural factors often outweigh headline narratives. Trump's victory, which prediction markets priced at approximately 60\% in the final week, was driven by a convergence of economic dissatisfaction, mobilization asymmetries, and a realignment of working-class voters\ldots''}
\end{quote}

These examples illustrate the core finding: all conditions produce fluent, professional-quality prose, but only the market-conditioned variants describe the correct outcome. The \nomarket article is arguably the most ``readable''---it tells a compelling story with rich detail---but it describes an event that did not occur.
